% subsubsection: 逆数型  \, $\displaystyle a_{n+1}=\frac{pa_n}{qa_n+r}$

\subsubsection{\texorpdfstring{逆数型  \, $\displaystyle a_{n+1}=\frac{pa_n}{qa_n+r}$}{逆数型}}
    では逆数型の漸化式について扱っていこう.
    分子が$a_n$の定数倍であることに注目する.
    逆数を取れば,分母にあった一次式を$a_n$で割り,定数と$\frac1{a_n}$の和に分けられる.
    以下では,$p\ne0$かつ$a_n\ne0$であり,$qa_n+r\ne0$として議論する.
    実際に見せよう.
    \begin{align*}
      &a_{n+1}=\frac{pa_n}{qa_n+r}
    \end{align*}
      両辺の逆数を取って
    \begin{align*}
      &\frac{1}{a_{n+1}}=\frac{qa_n+r}{pa_n}\\
      &\frac{1}{a_{n+1}}=\frac{q}{p}+\frac{r}{p}\cdot\frac{1}{a_n}
    \end{align*}
    $\displaystyle \frac{1}{a_n}=b_n$と置くと
    \begin{align*}
      &b_{n+1}=\frac{r}{p}b_n+\frac{q}{p}
    \end{align*}
    このようにすることで,特性方程式型の漸化式に帰着できる.

    \noindent \bookblocklink{example-reciprocal-1}{problem}{answer}{【例題】}
    \begin{enumerate}[label=(\arabic*),parsep=3pt,format=\bookitemlink{example-reciprocal-1}{problem}{answer}]
      \item $\displaystyle a_1=1,\, a_{n+1}=\frac{a_n}{a_n+1}$
      \item $\displaystyle a_1=1,\, a_{n+1}=\frac{-a_n}{2a_n+1}$
    \end{enumerate}

    \noindent\bookblocklink{example-reciprocal-1}{hint}{problem}{【方針】}
    \begin{enumerate}[label=(\arabic*),parsep=3pt,format=\bookitemlink{example-reciprocal-1}{hint}{problem}]
      \item とにかく基本に忠実に解きましょう.
      \item 逆数を取って一次の漸化式に帰着させましょう.
    \end{enumerate}

    \noindent\bookblocklink{example-reciprocal-1}{answer}{problem}{【解答】}
    \begin{enumerate}[label=(\arabic*),parsep=3pt,format=\bookitemlink{example-reciprocal-1}{answer}{problem}]
      \item $\displaystyle a_n=\frac{1}{n}$
      \item $\displaystyle a_n=\frac{1}{2(-1)^{n-1}-1}$
    \end{enumerate}

    \noindent\bookblocklink{example-reciprocal-1}{solution}{problem}{【解法】}
    \begin{enumerate}[label=(\arabic*),parsep=3pt,format=\bookitemlink{example-reciprocal-1}{solution}{problem}]
      \item $\displaystyle a_1=1,\, a_{n+1}=\frac{a_n}{a_n+1}$\\
            両辺の逆数を取って,
            \begin{align*}
              &\frac{1}{a_{n+1}}=\frac{a_n+1}{a_n}=1+\frac{1}{a_n}
            \end{align*}
            $b_n=\dfrac{1}{a_n}$と置くと
            \begin{align*}
              &b_{n+1}=b_n+1,\quad b_1=\frac{1}{a_1}=1\\
              &\therefore \quad b_n=1+(n-1)\cdot1=n\\
              &\therefore \quad a_n=\frac{1}{b_n}=\frac{1}{n}
            \end{align*}
      \item $\displaystyle a_1=1,\, a_{n+1}=\frac{-a_n}{2a_n+1}$\\
            両辺の逆数を取って,
            \begin{align*}
              &\frac{1}{a_{n+1}}=\frac{2a_n+1}{-a_n}=-2-\frac{1}{a_n}
            \end{align*}
            $b_n=\dfrac{1}{a_n}$と置くと,
            \begin{align*}
              &b_{n+1}=-b_n-2,\quad b_1=1
            \end{align*}
            特性方程式型として$\alpha$を求めると,
            \[
              \alpha=-\alpha-2 \ \Leftrightarrow\ \alpha=-1
            \]
            なので,\,$c_n=b_n+1$と置くと,
            \begin{align*}
              &c_{n+1}=-c_n,\quad c_1=b_1+1=2\\
              &\therefore \quad c_n=2\cdot(-1)^{n-1}\\
              &\therefore \quad b_n=c_n-1=2(-1)^{n-1}-1\\
              &\therefore \quad a_n=\frac{1}{b_n}=\frac{1}{2(-1)^{n-1}-1}
            \end{align*}
    \end{enumerate}
